\chapter*{Conclusiones}\label{chapter:Conclusions}
\addcontentsline{toc}{chapter}{Conclusiones}

En este estudio, se ha descrito el proceso de evaluación automática de soluciones vecinas en un Problema de Enrutamiento de Vehículos utilizando un grafo de evaluación. Este enfoque permite evaluar eficientemente las soluciones vecinas sin la necesidad de escribir código específico para esta tarea. 

A partir de este grafo, es posible obtener el valor de cualquier solución vecina al evaluarla en la función objetivo. Se ha resuelto una limitación de la evaluación automática, la cual impedía su aplicación en problemas que involucren penalizaciones durante el recorrido de una ruta.

Se agregó la cinta computacional que extiende este proceso a cualquier problema de enrutamiento de vehículos. Esta estructura permite conocer la parte de la solución afectada al aplicarle a la solución inicial un criterio de vecindad.

Por último, se ha descrito cómo el método de evaluación automática puede ser extendido a otras variantes del Problema de Enrutamiento de Vehículos, utilizando nuevos criterios de vecindad aplicables a cada una de ellas.